\section{Introduction}

Infectious disease outbreaks are commonly observed through univariate or multivariate time series of reported cases, such as the number of newly reported cases aggregated over a specified time interval. These data record when and how many cases were reported, but they do not directly reveal the underlying process by which infection spreads through a population. Mechanistic epidemic models connect the observed cases to changes in susceptible, infectious, and recovered populations by explicitly representing the dynamics of disease transmission \cite{breto2009mechanistic}. Partially observed Markov process (POMP) models provide a general framework for such models by separating two components: a latent Markov process describing the transmission dynamics and an observation process linking the latent states to the reported cases \cite{king2016pomp}.

Statistical inference for stochastic POMP models is generally difficult. Even when the static parameters and initial conditions are fixed, a stochastic model can generate many different latent epidemic trajectories. Evaluating the likelihood of the observed data requires integrating over these unobserved epidemic states and transmission histories, and the resulting likelihood therefore does not usually have a simple closed-form expression. Many statistical methods have been proposed for inference on POMP models \cite{breto2009mechanistic,king2016pomp}. These methods make it possible to use reported cases to infer epidemic processes and parameters that cannot be observed directly.

Within this framework, this report focuses on how to represent and infer a transmission rate that varies over time. The constant-$B$ POMP model assumes that the transmission rate remains fixed throughout the epidemic. In practice, the transmission rate may change with contact behavior, public-health interventions, and other conditions, while the transmission rate itself cannot be observed directly from reported cases \cite{bouman2024timevarying}. This report therefore incorporates Gamma noise into the process model and represents the transmission rate as a positive latent stochastic process \cite{breto2009mechanistic}. The Gamma-noise POMP model is then compared with the constant-$B$ POMP model to assess their ability to infer a time-varying transmission rate from reported cases.

This study is a simulation-based model comparison and a transmission-rate recovery study. Its central question is whether, under a prescribed piecewise truth in which the transmission rate decreases from 4 to 2 beginning at week 5, the Gamma-noise POMP model can recover the known time-varying transmission rate more accurately than the constant-$B$ POMP model. The experimental design contains three distinct POMP models. A data-generating POMP combines the prescribed transmission-rate path with a stochastic SIR transmission process and a negative-binomial observation model to generate reported cases. The same reported cases are then supplied to two fitting POMPs. The Gamma-noise POMP model treats the transmission rate as a time-varying latent stochastic process, whereas the constant-$B$ POMP model estimates a single transmission-rate parameter for the entire epidemic. Neither fitting model is given the true transmission-rate path or the location of its change point.

For a fixed set of parameter values, the particle filter is used to estimate the likelihood of the observed data and obtain filtering summaries of the transmission rate and other latent states \cite{gordon1993particle,king2016pomp}. Iterated filtering (IF2) uses this particle-filtering framework to search the static parameter space of each fitting model \cite{ionides2015if2}. These algorithms support statistical inference for the model comparison; they are not components of the epidemic models themselves. The inferential performance of the two models is evaluated primarily by comparing the estimated transmission rates with the known truth using recovery error and root mean squared error (RMSE).

This report investigates the ability of the Gamma-noise POMP model to infer a time-varying transmission rate and compares its performance with that of the constant-$B$ POMP model in a controlled simulation study. The results show that, when the prescribed transmission rate decreases from 4 to 2, the Gamma-noise POMP model recovers the transmission-rate path more accurately and produces a lower RMSE. Section~2 defines the stochastic epidemic process, the observation process, and the two transmission models. Section~3 introduces the particle filter and IF2 used for inference, together with the measures used to evaluate recovery performance. Sections~4--6 present the simulation design, the results, and their interpretation.
